% ARCHIVED at v3.88 (2026-09-24) from chapters/06_results.tex -- author: "6.3.2.4 Conclusion, too long, is the model no diff? only K matters? ... did the HF beat the PCC? that is the RQ, plz answer clearly". Replaced by a two-paragraph conclusion answering RQ1 and RQ2.
% Kept verbatim for rollback; not part of the build.

\subsubsection{Conclusion}\label{sec:res:uav:corridor:conclusion}

UAV-corridor asks whether the projector repairs a plan against a constraint that leaves the horizontal plane
and that no demonstration satisfies. It does so for what the constraint is about. Projected, the vehicle
descends under the leaned plane on $540$ of $560$ flights, by a median of $0.12$ to $0.37$\,m, and climbs
over the roof on all $560$, by $0.27$ to $0.39$\,m, where the unprojected plans fly level through both; and
the violating steps of a flight fall from $74.5$ to $99.4$ under the tilt and $28.6$ to $36.6$ over the hump
to $1.4$ to $10.4$ and $0.2$ to $8.3$ for the flow-based models. It does not remove every violation. What
remains lies where the constraint ends: under the tilt in the last $0.3$\,m of the leaned plane, where the
commanded position has already left the plane's span and the plan is released from it while the vehicle,
half a metre behind, is still inside; over the hump at its apex. The budget is the trade that shrinks this
residue --- every step up in network evaluations buys fewer violating steps with more time per action --- and
at twenty evaluations it is down to one or two steps per flight under the tilt and a fifth of one over the
hump. At the other end of the ladder, per-step projection at one and two evaluations cannot carry a plan over
the hump: the plans stay within the altitude range of the demonstrations while the apex asks for $1.48$\,m,
and the flights stop short of it. The diffusion baseline leaves the fewest violating steps, $0.8$ and $0.0$,
and is the slowest planner of the scene, at $654$ to $666$\,ms per action and $290$ and $321$ control steps to
the end of the corridor, against $247$ to $316$ for the flow-based configurations.

The central finding of the scene concerns the models: the generative model does not matter on it, and the
reason lies in its demonstrations. They are straight lines: each runs along the corridor from $x=-2.8$ to
$2.8$\,m on one of three lanes, $y\in\{-0.12,0,0.12\}$\,m, at an altitude drawn once per episode from
$[0.90,1.30]$\,m (\autoref{tab:uav-demos}, \autoref{fig:expert-uav}), so that from any observed position the demonstrations
continue along a single path. Before projection the flights of the four models cannot be told apart: every
flight reaches the end of the corridor, in $242$ to $251$ control steps on average, and the violating steps
of any two models lie within about one standard deviation of the flights. After projection the three
flow-based objectives, wherever their flights reach the end of the corridor, leave the same violating steps
at a shared budget and projection method to within $1.6$ per flight, at most two standard deviations of the
flights, while under the tilt the budget takes them from about ten at one evaluation to one or two at
twenty; in control steps to the finish line they differ by at most fifteen of about $250$ to $300$, and no
objective is the fastest at every budget. What orders the flow-based configurations of UAV-corridor is the
budget and the projection method, not the objective. This is the case the average-velocity objectives are
not built for: they correct the curvature that arises where the straight per-sample paths of many different
continuations are averaged into one field (\autoref{sec:bg:fewstep}), and where the demonstrations continue
along a single path there is little to average, so that at one evaluation instantaneous-velocity matching
already arrives where they do.
\guard{One training seed per model. The curvature of the learned fields is not measured; what the scene
measures is the tie.}

The one model that still differs after projection is the diffusion baseline, and against it the relationship
D3IL-avoiding established between the models and the projection holds in part
(\autoref{sec:res:avoiding:conclusion}). What carries over is the price: the flow-based models reach the end
of the corridor on every flight at a small fraction of the baseline's time per action, a twenty-fourth under
the tilt and a ninth over the hump at their cheapest, and at equal violation they reach it in fewer control
steps than the baseline in the group that contains it. What does not carry over is the removal of every
violation: on D3IL-avoiding the flow-based models at one and two evaluations succeed within the constraints
on almost every episode, while here their flights keep a residue of violating steps at every budget, and the
baseline leaves the fewest. Between the projection methods the budget decides: at three and five
evaluations endpoint projection is as cheap or cheaper and leaves as many violating steps or more, and at
twenty it leaves fewer and shallower ones under both constraints; over the hump with a single candidate it
overtakes the per-step projector at a lower cost, $0.2$ against $1.8$ violating steps per flight at $291$
against $355$\,ms, the bar this thesis set for it.
\dataref{v3.86 (author: \enquote{where is the conclusion in the UAV-corridor section?}); v3.87 (author:
\enquote{in the expert data very simple case the Model NOT matter, this is KEY CONCLUSION, our expert data for
corridor just a line, quote the chap5}). Demonstrations: \autoref{tab:uav-demos} (\texttt{trajectories.py}:
\texttt{corridor\_path}, \texttt{CORRIDOR\_CHANNELS}; \texttt{generator.py}: \texttt{\_build\_traj\_and\_init}).
The tie, MeanFM / CI-MeanFM / FM: unprojected violating steps tilt $74.5$--$99.4$ (sd $25.6$--$34.8$), hump
$28.6$--$36.6$ (sd $6.8$--$8.9$), all four models; projected, the cells that reach the end of the corridor:
tilt per-step $\nfe=1$ $10.2/10.3/10.2$, $\nfe=2$ $10.3/10.4/9.8$, $\nfe=3$ $6.6/5.0/5.8$ (the $1.6$; sd $0.8$),
endpoint $\nfe=3$ $7.6/7.1/6.9$; hump per-step $\nfe=3$ $4.9/4.6/5.6$, endpoint $\nfe=3$ $8.0/8.0/8.3$; the
largest gap in steps, hump endpoint $\nfe=3$, $284.6$ against $299.3$. The rest are the facts of
\autoref{sec:res:uav:corridor:raw}--\autoref{sec:res:uav:projection}, nothing new: altitude and the hand-over
(\texttt{DA\_20260924\_corridor\_v3.md} \S3.6, \S4.1, re-verified at v3.80: commanded position $0.49$--$0.61$\,m
ahead at the first violating step, clear of the constraint at $4014/4014$ violating steps), the stop at
$\nfe\le2$ (\S4.2), the end-of-corridor line (R45a), the tables and the frontier
(\texttt{data/corridor\_v3\_frontier.json}). Price: $666.3/27.4 = 24$ (tilt, MeanFM $\nfe=1$), $654.2/71.9 = 9.1$
(hump, CI-MeanFM $\nfe=3$ endpoint). D3IL-avoiding: \autoref{tab:avoiding-dpcc-protocol} (flow models at
$\nfe=1,2$, S\&C $0.90$--$1.00$). v3.87: over the hump at $\nfe=5$ endpoint and per-step projection leave the
same $3.9$, hence \enquote{as many or more}.}
